\section{How does temporal predictability influence neural and behavioral responses to repeated sounds?}
\label{sec:q7}

% Part III — Updating responses over time (Q7).
% See docs/OUTLINE.md — Q7.
% Draft 1. Style models: Watabe-Uchida + Blumstein (agents/STYLE.md).

% VERIFY: LimaBednekoff1999 (not full-text checked): wording of the risk allocation predictions.
The timing of repeated sounds carries information independently of their acoustic identity. At an ecological timescale, the predation risk allocation hypothesis (\secref{sec:intro-framework}) predicts that the same alarm call should evoke less vigilance in a receiver that has heard alarms often than in one that has heard them rarely, for reasons unrelated to sensory habituation or caller reliability \citep{LimaBednekoff1999}. Because it has not been tested with acoustic cues in rodents, it motivates the question here rather than answering it.

At shorter timescales, the relevant concept is temporal expectation. A useful description of a schedule is its hazard function: the probability that a sound will occur in the next moment, given that it has not occurred since the last one. On a fixed schedule the hazard rises sharply as the expected time approaches, so the receiver can predict when the next sound will occur. On a schedule in which intervals follow an exponential distribution, the hazard is constant, so elapsed time provides no information about when the next sound will arrive. A receiver that habituates partly by predicting the timing of stimulation should therefore habituate faster to regular than to irregular schedules, and most slowly to schedules with constant hazard. This is a hypothesis derived from the concept, not a demonstrated optimization rule, and testing it requires separating the variability of intervals from their mean length, from the time available for recovery between sounds, and from total exposure.

Laboratory studies with artificial sounds support the first part of this prediction. In rats, regularly timed startle stimuli produced more habituation than irregularly timed stimuli at the same 8-s mean interval \citep{Davis1970}. In the same study, 2-s spacing produced more decline during training than 16-s spacing, but 16-s spacing produced lower startle when animals were retested, so interval length affects short-term and long-term habituation differently and must be separated from regularity. In mice, temporally jittered trains of neutral tone pulses prevented the rapid habituation of basolateral amygdala responses that regular trains produced, and the jittered trains led to avoidance \citep{HerryEtAl2007}. Herry and colleagues used two regimes: 5-Hz trains with an interval standard deviation of 140~ms for behavior and c-Fos, and 0.7-Hz trains with a mean interval of 1429~ms and a standard deviation of 1368~ms for electrophysiology. Both studies manipulated the variability of intervals, not the shape of their distribution, and both used artificial stimuli: a startle tone and a neutral tone. They show that irregular timing can preserve responsiveness, but they do not test the hazard-based prediction or whether biologically meaningful sounds interact differently with temporal prediction.

Dopamine recordings indicate that prediction changes how accumbens dopamine recovers after an aversive sound. In freely moving rats, accumbens core dopamine fell during aversive noise and during a cue that predicted it, and recovery after unexpected noise was slower than after predicted noise \citep{GoedhoopEtAl2022}.
% VERIFY: KutluEtAl2025 (not full-text checked): species, stimuli, and time course of novelty effects on accumbens dopamine.
Novelty also altered accumbens dopamine responses to aversive stimuli over windows extending beyond the novel event \citep{KutluEtAl2025}.
These results show that a neural measure of response to an aversive sound depends on what the animal expected, which is the minimal condition for timing to influence responses through prediction. They do not show that dopamine encodes the hazard of a sound or that it mediates the behavioral effects of irregular schedules.

Field studies show that free-living rodents decode the temporal structure of natural calls, but not uniformly.
% VERIFY: BlumsteinArmitage1997, SloanEtAl2005, SwanHare2008, SloanHare2006 (not full-text checked): temporal structure and rate effects on vigilance.
Marmots use call rate as information about risk \citep{BlumsteinArmitage1997}, and in Richardson's ground squirrels repetition and call structure alter vigilance \citep{SloanEtAl2005,SwanHare2008}. Yet adult Richardson's ground squirrels ignored experimental changes in the rate and bout length of juvenile calls, although juveniles themselves use call rate as a cue to urgency \citep{SloanHare2006}. Temporal structure within a calling bout is thus informative to some receivers and not others. These studies concern the temporal pattern of a single bout, not the distribution of intervals between bouts across hours or days, which is the variable that risk allocation and temporal expectation address.

The empirical conclusion is therefore narrow. Regular timing can promote habituation and irregular timing can preserve responsiveness in some laboratory paradigms with artificial sounds. Field work has examined call rate, repetition, sequence, caller identity, and reliability, but we found no study that compared interval distributions (fixed, uniform, exponential, or heavy-tailed) using natural rodent calls, and constant-hazard timing remains a hypothesis rather than a result. \secref{sec:conclusion} outlines a direct test that holds event count, mean interval, and acoustic exposure constant across schedules differing only in hazard. Whether any temporal statistic that preserves responsiveness to artificial sounds also preserves it to natural calls is an open question.

\tabref{tab:summary-q7} consolidates these claims for later condensation or figure drafting.

% Intermediate summary matrix for Q7. May later be condensed, moved to an appendix, or redrawn as a figure.
\begin{table}[tbp]
\centering
\footnotesize
\singlespacing
\caption{Working summary for Q7: temporal predictability. Hazard and risk allocation are theory; lab timing effects use artificial sounds; field work addresses within-bout structure, not interval distributions.}
\label{tab:summary-q7}
\begin{tabular}{@{}p{1.35in}p{2.55in}p{1.2in}p{0.9in}@{}}
\toprule
\textbf{Factor} & \textbf{Finding} & \textbf{Stimulus class} & \textbf{Status} \\
\midrule
Risk allocation & Less vigilance per alarm when high-risk periods are frequent & Theory & Untested acoustically \\
\addlinespace[2pt]
Regular vs irregular ISI & Regular timing can promote habituation more than irregular at matched mean & Startle / tones & Strong (artificial) \\
\addlinespace[2pt]
Mean interval length & Short- and long-term habituation respond differently & Startle tones & Supporting \\
\addlinespace[2pt]
Hazard / distribution shape & Constant-hazard (exponential) predicted to preserve responding most & Hypothesis & Open \\
\addlinespace[2pt]
Prediction (DA recovery) & Unexpected noise slows NAc dopamine recovery & Aversive noise & Mechanistic bridge \\
\addlinespace[2pt]
Within-bout rate / structure & Call rate and repetition alter vigilance; age-dependent nulls exist & Natural sciurid calls & Strong (bout-level) \\
\addlinespace[2pt]
Natural calls $\times$ schedule & Interval distributions compared with natural calls & --- & Gap \\
\bottomrule
\end{tabular}
\end{table}

